% TOP_PROBLEM: {"id": 30007160, "title": "Spectral scalar-curvature control of width in dimension three", "category_id": 6, "category": {"id": 6, "name": "geometry", "display_name": "Geometry", "description": "Euclidean and non-Euclidean geometry, geometric structures.", "slug": "geometry", "order_index": 6, "created_at": "2026-07-31T15:26:25.671Z"}, "status": "open"}
% FIRST_POSTED: 2026-10-06
% ATTEMPT_STATUS: solved
% OPENAI_CLAIM: {"schema_version":1,"source_repo":"https://github.com/openai/math","source_commit":"adc7f1241b42e322a6451854ab7e4b4c146bf78a","release_date":"2026-10-06","match":"full","resolution":"proved","independently_reviewed":false,"second_pass":"agreed","families":[{"family":"336","title":"Spectral scalar curvature, Urysohn width, and macroscopic dimension","subject":"Differential geometry","manuscripts":[{"title":"Spectral scalar curvature and Urysohn width in dimension three","pdf_path":"preprints/Spectral-scalar-curvature-and-Urysohn-width-in-dimension-three-October-5-2026/spectral-urysohn-three-manifolds.pdf","readme_path":"preprints/Spectral-scalar-curvature-and-Urysohn-width-in-dimension-three-October-5-2026/README.md","date":"2026-10-05","theorem_ref":"preprints/Spectral-scalar-curvature-and-Urysohn-width-in-dimension-three-October-5-2026/build/sections/01-introduction.tex#thm:main"}],"lean":null,"reasoning_trace":null}]}
\documentclass[11pt]{article}
\usepackage[margin=25mm]{geometry}
\usepackage{fontspec}
\setmainfont{Latin Modern Roman}
\usepackage{amsmath,amssymb,mathtools}
\usepackage{unicode-math}
\setmathfont{Latin Modern Math}
\usepackage{enumitem}
\usepackage{xurl,hyperref,bookmark}
\hypersetup{colorlinks=true,urlcolor=blue}
\setcounter{secnumdepth}{0}
\setlength{\parindent}{0pt}
\setlength{\parskip}{5pt}
\setlength{\emergencystretch}{3em}
\begin{document}
\section*{Spectral scalar-curvature control of width in dimension three}

\subsection{OpenAI's claim}
\textit{Spectral width in dimension three}: A complete three-manifold satisfying $-4\Delta+\mathrm{Scal}\ge\lambda>0$ is claimed to map continuously to a graph with full fiber diameter at most $500/\sqrt\lambda$. No global bounded-geometry condition is assumed.

\subsection{Outline of the argument}
\textit{Spectral width in dimension three}: A positive supersolution is converted into a drift-weighted scalar inequality without changing the original distance. Weighted stable surfaces inherit a two-dimensional inequality, whose geodesic index estimate bounds their intrinsic diameter. Local minimization and smoothing in compact regions then organize these surfaces into a graph-valued decomposition with a uniform bound on each complete fiber.

\subsection{Scope relative to this problem}
\textit{Spectral width in dimension three}: A complete three-manifold satisfying $-4\Delta+\mathrm{Scal}\ge\lambda>0$ is claimed to map continuously to a graph with full fiber diameter at most $500/\sqrt\lambda$. No global bounded-geometry condition is assumed. The parameters, regularity assumptions and normalization in this question are retained in applying that assertion.

With the hypotheses and conventions above, the claimed result addresses the whole mathematical assertion on this page.

\subsection{Formal verification}
\textit{Spectral width in dimension three}: The pinned catalogue lists no Lean formalization for this result.

This site has not run the Lean development. These links describe OpenAI’s selected formal statements and should be read with their stated scope.

\subsection{Status caveat}
OpenAI reports that an internal model produced the manuscripts. Their arguments have not been independently verified as proofs by this site. Any portion outside the documented formal scope may contain mistakes, and the catalogue label reports the announced mathematical claim.

\subsection{Sources}
\begin{itemize}
\item OpenAI, \emph{Spectral scalar curvature and Urysohn width in dimension three}, 2026-10-05; citation key \texttt{OAI:Spectral-scalar-curvature-and-Urysohn-width-in-dimension-three-October-5-2026}. \href{https://github.com/openai/math/blob/adc7f1241b42e322a6451854ab7e4b4c146bf78a/preprints/Spectral-scalar-curvature-and-Urysohn-width-in-dimension-three-October-5-2026/spectral-urysohn-three-manifolds.pdf}{PDF (GitHub)}; \href{https://github.com/openai/math/raw/adc7f1241b42e322a6451854ab7e4b4c146bf78a/preprints/Spectral-scalar-curvature-and-Urysohn-width-in-dimension-three-October-5-2026/spectral-urysohn-three-manifolds.pdf}{PDF (direct)}; \href{https://github.com/openai/math/blob/adc7f1241b42e322a6451854ab7e4b4c146bf78a/preprints/Spectral-scalar-curvature-and-Urysohn-width-in-dimension-three-October-5-2026/README.md}{Manuscript page and citation}; \href{https://github.com/openai/math/blob/main/preprints/Spectral-scalar-curvature-and-Urysohn-width-in-dimension-three-October-5-2026/spectral-urysohn-three-manifolds.pdf}{Latest version}.
\item \href{https://github.com/openai/math/blob/adc7f1241b42e322a6451854ab7e4b4c146bf78a/CONTENTS.md}{OpenAI catalogue entry}.
\end{itemize}
\end{document}
